\subsection{3 强化学习攀登}\label{the-reinforcement-learning-climb}

预训练和中训练使 MAI-Base-1 具备了广泛的预测能力与知识，但它们并未指定模型应如何行事、如何解决长时程任务，或如何分配推理时计算。我们的强化学习（RL）攀登通过针对任务特定反馈优化模型来解决这一问题，使模型能够在作答前生成思维链（CoT），使用外部工具，与环境交互，并遵循偏好与安全信号。

由于 MAI-Thinking-1 是我们首款自研推理模型，我们的 RL 攀登从一个从未接触过推理轨迹的检查点开始。因此，模型必须从零开始发展推理能力，这使得长期训练稳定性成为核心挑战。我们通过三种机制实现这一点：（i）对组相对策略优化（GRPO，Shao et al., 2024）进行两项简单但关键的调整，（ii）在崩溃或基础策略更新后通过自蒸馏恢复 RL 攀登（第 3.1 节），以及（iii）消除训练与推理之间数值失配的基础设施改进（第 3.6 节）。

为支持并行开发，我们训练三个领域特定的专家模型：一个面向 STEM 与竞技编程（第 3.2 节），一个面向智能体编码与工具使用（第 3.3 节），一个面向有用性与安全性（第 3.4 节）。所有专家的 RL 攀登都遵循相同的配方，但在提示词分布以及模型因何获得奖励方面有所不同。如图 12 所示，这些专家模型随后通过监督微调（SFT）整合为单一模型。最后一轮轻量级 RL 攀登将这一整合模型转化为 MAI-Thinking-1——一个在所有领域都表现强劲的模型（第 3.5 节）。



\pandocbounded{\includegraphics[keepaspectratio,alt={image}]{images/4386f66d5e807c10b3d62b8f9332b4d8a6f2fc5756b594189318ecfcb85254b5.jpg}}

图 12. 我们各 RL 攀登的概览。我们从 MAI-Base-1 出发，使用强化学习训练三个专家模型。这些专家随后通过简单 SFT 被蒸馏为单个整合模型，该模型再经历一轮最终 RL 阶段，产生 MAI-Thinking-1。

\subsection{3.1 强化学习配方}\label{reinforcement-learning-recipe}

\subsection{3.1.1 RL 目标}\label{rl-objective}

我们的强化学习攀登从一个策略 \(\pi _ { \theta }\)（即我们的中训练模型）出发。对于提示词 \(q ,\)，rollout 策略采样一组共 G 个响应 \(y _ { 1 : G }\)，每个响应 \(y _ { i }\) 会收到一个标量奖励：

\[
R _ { i } = R ( q , y _ { i } ) .\tag{4}
\]

奖励函数 R 与领域相关；通常它要么基于代码执行，要么基于提示式 AI 裁判或训练好的奖励模型的反馈。我们从 GRPO（Shao et al., 2024）出发推导训练目标，并采用 token 级策略梯度（Yu et al., 2025）：

\[
J ( \theta ) = \mathbb { E } _ { q \sim P ( Q ) , y _ { 1 : G } \sim \pi _ { \mathrm { o l d } } } \left[ \frac { 1 } { \sum _ { i = 1 } ^ { G } | y _ { i } | } \sum _ { i = 1 } ^ { G } \sum _ { t = 1 } ^ { | y _ { i } | } \operatorname* { m i n } \Bigl ( r _ { i , t } ( \theta ) A _ { i } , \ \mathrm { c l i p } \bigl ( r _ { i , t } ( \theta ) , 1 - \epsilon , 1 + \epsilon \bigr ) A _ { i } \Bigr ) \right] ,\tag{5}
\]

其中 \(P ( Q )\) 是所有提示词上的分布，\(\pi _ { \mathrm { o l d } }\) 表示用于生成 rollouts 的策略。实践中，归一化在全局训练批次（跨所有数据并行 rank）的所有 token 上计算，使每个 token 无论响应长度如何都贡献相同。对于响应 \(y _ { i }\) 和 token 位置 \(t ,\)，重要性采样比率为

\[
r _ { i , t } ( \theta ) = \frac { \pi _ { \theta } ( y _ { i , t } \mid q , y _ { i , < t } ) } { \pi _ { \mathrm { o l d } } ( y _ { i , t } \mid q , y _ { i , < t } ) } .\tag{6}
\]

响应级优势
\(A _ { i } = \left( R _ { i } - \operatorname* { m e a n } ( R _ { 1 : G } ) \right) / \operatorname { s t d } ( R _ { 1 : G } )\)
在响应 \(y _ { i }\) 的所有 token 之间共享。

我们对该目标应用两项修改：自适应熵控制，动态调整上裁剪界以维持目标策略熵；以及外层比率裁剪，对我们的目标中未裁剪的分支设置上限，以防止梯度范数爆炸。

自适应熵控制。与 Yu et al.~(2025) 和 Mistral-AI et al.~(2025) 类似，我们使用分离的下、上裁剪界。我们用单个基础超参数 ϵ 对它们进行参数化，它控制基础信赖域宽度，再加上与熵相关的对裁剪上界的松弛量 k：

\[
r _ { i , t } ^ { \mathrm { t r } } ( \theta ) = \mathrm { c l i p } \left( r _ { i , t } ( \theta ) , 1 - \epsilon , ( 1 - \epsilon ) ^ { - 1 } + k \right) .
\]



图 13. 训练期间的自适应熵控制。上图：目标熵 \(H ^ { \star } = 0 . 3\) 的一次 RL 攀登中，800 步内观测到的熵 \(\hat { H } ( \pi _ { \theta } )\)。下图：\(k\) 的对应取值，其中
\(k _ { \operatorname* { m a x } } = 1 . 0 \colon\) 当观测熵超过目标时减小，否则增大。
调整 k 为调节策略熵提供了有效机制。

\pandocbounded{\includegraphics[keepaspectratio,alt={image}]{images/e8cf2dc9009381a2f7dc7672e0bf0e8e8be51f7785c9c0e1e996878378daa61c.jpg}}

我们发现上界需要仔细调参：过大时会导致熵爆炸，过小时会导致熵坍缩（另见 Cui et al., 2025; Xi et al., 2025）。

为解决这一问题，我们基于当前策略的熵，使用一个简单的积分控制器动态调整 k。在每个训练步，我们通过重要性加权估计器估计目标策略的每 token 熵：

\[
\hat { H } ( \pi _ { \theta } ) = \frac { 1 } { | T | } \sum _ { ( i , t ) \in \mathcal { T } } - \log \pi _ { \theta } ( y _ { i , t } \mid q , y _ { i , < t } ) \cdot r _ { i , t } ( \theta ) ,\tag{7}
\]

其中 \(\mathcal { T }\) 表示当前训练批次中所有（响应索引，token 索引）对的集合。给定目标熵 \(H ^ { \star }\)，控制器在每步之后以步长 \(\delta \in \mathbb { R } ^ { + }\) 更新 k：

\[
k \gets \mathrm { c l i p } \left( k + \delta \cdot \mathrm { s i g n } \left( H ^ { \star } - \hat { H } ( \pi _ { \theta } ) \right) , 0 , k _ { \mathrm { m a x } } \right) .\tag{8}
\]

直观上，当熵过低时，增大 k 会放宽上裁剪界，使策略能够更激进地提高备选 token 的概率。当熵足够高时，减小 k 以收紧信赖域。如图 13 所示，该机制起到自动熵正则化器的作用，而无需在损失中添加显式的熵奖励项——我们发现后者不如自适应熵机制。

我们将 \(k = 0\) 作为初值，使裁剪界 1 − ϵ 与 \(( 1 - \epsilon ) ^ { - 1 }\) 互为乘法逆元，从而让初始裁剪区间在对数比率空间上对称。随后我们按照上述熵控制器在线更新 k。

外层比率裁剪。GRPO 目标刻意保留两种情况不裁剪：（a）优势为负且新策略赋予比旧策略更高的概率（即 \(A _ { i } < 0\) 且 \(r _ { i , t } > 1\)），以及（b）优势为正且新策略赋予更低的概率（\(A _ { i } > 0\) 且 \(r _ { i , t } < 1\)）。其最初动机（Schulman et al., 2017）是：当策略朝正确方向自我修正时不加限制，仅限制利用优势估计的移动。



然而在实践中，我们发现这些未裁剪分支有时会导致灾难性的梯度范数尖峰。我们通过在所有分支上施加硬外层裁剪来解决这一问题：

\[
r _ { i , t } ^ { \mathrm { o u t } } ( \theta ) = \mathrm { c l i p } \big ( r _ { i , t } ( \theta ) , r _ { \mathrm { m i n } } , r _ { \mathrm { m a x } } \big ) ,\tag{9}
\]

其中 \(r _ { \mathrm { m a x } }\) 设为较大的值，而 \(r _ { \mathrm { m i n } }\) 可以不加约束；这与 Ye et al.~(2020) 提出的双重裁剪 PPO 目标在精神上类似。

这一两层策略在保留正常范围内比率的标准信赖域行为的同时，丢弃新旧概率存在极端差异的情况。实验表明，这减少了梯度范数尖峰，使攀登更加稳定。

\subsection{3.1.2 奖励设计}\label{reward-design}

虽然我们的任务特定组件在不同领域的 RL 攀登中有所不同，但我们在整个过程中使用相同的奖励分解：

\[
R ( q , y _ { i } ) = R _ { \mathrm { t a s k } } ( q , y _ { i } ) + w _ { \mathrm { l a n g } } \cdot R _ { \mathrm { l a n g } } ( y _ { i } ) - w _ { \mathrm { l e n } } \cdot R _ { \mathrm { l e n } } ( y _ { i } ) .\tag{10}
\]

这里，\(R _ { \mathrm { t a s k } }\) 表示任务特定奖励，\(R _ { \mathrm { l a n g } }\) 是语言一致性奖励，\(R _ { \mathrm { l e n } }\) 是长度惩罚。系数 \(w _ { \mathrm { l a n g } }\) 和 \(w _ { \mathrm { l e n } }\) 是标量超参数。

语言一致性奖励。随着 RL 期间上下文长度的增加，我们观察到模型开始在其 CoT 中产生外语 token。这些混合语言 CoT 与训练和推理策略之间对数概率散度的尖峰相关，最终会破坏训练稳定性。为解决这一问题，我们使用与 Guo et al.~(2025) 类似的语言一致性奖励 \(R _ { \mathrm { l a n g } } ( y _ { i } )\)。由于英语在我们的训练分布中占主导地位，我们针对英语定义语言一致性：

\[
R _ { \mathrm { l a n g } } ( y _ { i } ) = \operatorname* { m a x } \Bigl ( 1 - \alpha \cdot n _ { \mathrm { n o n - e n g l i s h } } ( y _ { i } ) , 0 \Bigr ) ,\tag{11}
\]

其中 \(n _ { \mathrm { n o n - e n g l i s h } } ( y _ { i } )\) 是响应 \(y _ { i }\) 的 CoT 中非英语单词的数量，$\alpha$ 是逐词惩罚项。实践中我们发现，top-p 采样在阻止单个低概率外语 token 方面同样有效。

长度惩罚。沿用 Xiang et al.~(2025)，我们将长度惩罚定义为：

\[
R _ { \mathrm { l e n } } ( y _ { i } ) = \rho _ { q } \cdot \frac { | y _ { i } | } { \ell _ { \mathrm { m a x } } }\tag{12}
\]

其中 \(\rho _ { q }\) 是问题 q 的通过率，\(\ell _ { \mathrm { m a x } }\) 是最大 rollout 长度。该惩罚同时取决于响应长度和问题难度。通过率低的难题受到更弱的惩罚，允许模型探索更长的推理轨迹。简单题受到更强的惩罚，通过去除冗余循环和回避性行为，鼓励简洁且成本高效的推理。

\subsection{3.1.3 采样策略}\label{sampling-strategy}

我们采用多种采样策略来提高 RL 攀登的效率和稳定性。它们作用于两个层面：选择训练哪些问题，以及控制如何为这些问题生成 rollouts。



用户与助手之间的对话。用户提出问题，助手负责解答。助手先在头脑中思考推理过程，再向用户给出答案。推理过程和答案分别放在 <think> </think> 和 <answer> </answer> 标签中，即 <think> 此处为推理过程 </think> <answer> 此处为答案 </answer>。用户：\{QUERY\} 助手：<think> \\
图 14. 用于第一轮自蒸馏之前的推理提示词模板，其中 \{QUERY\} 是实际用户查询的占位符（来自 Guo et al.~(2025)）

问题采样。对于训练集中的每个问题 q，我们从当前推理模型生成一组 rollouts。为降低推理成本，我们采用提前退出策略：先采样 \(G _ { \mathrm { e a r l y } } \ < \ G\) 个响应并计算其经验通过率，即奖励为正的响应所占比例。如果早期通过率落在可接受区间 \([ \rho _ { \mathrm { m i n } } ^ { \mathrm { e a r l y } } , \rho _ { \mathrm { m a x } } ^ { \mathrm { e a r l y } } ]\) 内，则采样全部 G 个响应；否则丢弃该问题。

在生成全部 G 个响应之后，我们对完整组应用第二次通过率过滤 \([ \rho _ { \mathrm { m i n } } , \rho _ { \mathrm { m a x } } ]\)。只有完整通过率落在该范围内的问题才会用于训练。这第二次过滤的主要动机是剔除方差极低的分组：如果几乎全部响应都正确或都不正确，该分组能提供的相对学习信号很有限。

Rollout 采样。我们使用 top-p 采样（Holtzman et al., 2019），用 \(\pi _ { \mathrm { o l d } }\) 采样 rollouts \(y _ { 1 : G }\)。我们发现，继续对采样核之外 token 对应的 logits 进行反向传播可能导致灾难性的 off-policy 失配，使训练在几步之内就发散。为防止这一点，我们遵循 DeepSeek-AI et al.~(2025b) 的做法，在训练期间排除这些 token：复用 rollout 采样时的 top-p 截断掩码，在 softmax 计算之前将所有被排除 token 的 logits 设为 −$\infty$。我们发现 top-p 掩码显著降低了 RL 训练期间的策略发散，代价是掩码存储与重放带来的额外开销。

为了提高攀登早期阶段的训练效率，我们最初将最大 rollout 长度限制在 8k token。随着训练推进，我们以 2 的幂逐级提高该上限，直到达到 128k token 的最终输出长度。我们发现这一长度扩展课程能显著降低低性能阶段的推理成本——该阶段很少需要长推理轨迹——同时仍能让模型随能力提升逐步适应更长的上下文。

\subsection{3.1.4 自蒸馏}\label{selfdistillation}

从中训练检查点实现强性能需要大量 RL 步的训练。我们利用自蒸馏（Huang et al., 2022; Zelikman et al., 2022; Singh et al., 2023; Zhao et al., 2026）让此类长时运行攀登更切实可行。自蒸馏的具体做法是：收集 RL 期间生成的 rollouts，并在一个中训练检查点上对这些 rollouts 做 SFT。得到的模型作为继续 RL 攀登的起点，同时保留先前 RL 阶段发现的能力。我们将自蒸馏用于多个目的：

• 对于最初的攀登，我们使用领域或任务特定的提示词来激发目标行为；例如，图 14 展示了我们 STEM 攀登使用的初始原始文本提示词。自蒸馏使我们只需更新 SFT 数据的格式，就能从提示词驱动迁移到使用原生对话格式。

• 另一个主要用途是从偶发的运行失败中恢复。在我们 RL 技术栈的早期——即第 3.6.4 节所述稳定性改进引入之前——推理与训练之间微小的数值差异有时会在训练过程中累积，导致攀登发散。在这种情况下，自蒸馏提供了一种简单有效的方式，将一次 RL 攀登的进度承接给下一次。从崩溃前的检查点恢复这一备选方案往往不可行，因为某些不稳定性在崩溃实际发生前许多步就已嵌入模型参数之中。



图 15. STEM 攀登期间在 AIME 2025（左）和 LiveCodeBench v6 硬子集（右）上的表现。自蒸馏通过 ⋆ 标记标出；不同的预训练/中训练模型版本以不同颜色显示。训练全程的最大输出长度在底部标出。自蒸馏是一种在崩溃后重置数值（可通过性能的突然下降观察到）以及在我们运行过程中更换基础策略的有效方式。随着我们在攀登中持续做出基础设施与算法层面的改进，因运行崩溃而触发的自蒸馏越来越少见。

\pandocbounded{\includegraphics[keepaspectratio,alt={image}]{images/426a11231cea81174de083d73fd1712f6827659c20fa8c0cca290ef4257fe3db.jpg}}

• 当新的预训练和中训练检查点可用时，自蒸馏使我们能够将先前攀登的进展延续到下一代模型。

• 在自蒸馏过程中，可以应用过滤器剔除任何表现出不当奖励劫持形式的样本。

图 15 展示了自蒸馏在 STEM 攀登中的两种用途：既用于从数值不稳定中恢复，也用于更新基础策略。

关键发现与最佳实践。我们进行了大量消融实验，以弄清在实践中如何最好地进行自蒸馏。主要发现如下：

• O(1M) 量级的推理轨迹足以匹配教师性能，同时保留 SFT 的稳定性优势。使用大得多的数据集收益递减，且有过度约束策略的风险——它会收窄策略的输出分布，使 RL 恢复后几乎没有探索空间（Zheng et al., 2024; Gao et al., 2025）。

• 使用包含导致最终答案错误的轨迹进行训练，与仅使用成功轨迹训练效果相当（Gandhi et al., 2025）。由于我们的 RL 运行通常产生的成功轨迹远超 O(1M)，我们最终将训练限定在成功轨迹上。



• 使用攀登后期阶段的轨迹很重要。纳入非常早期检查点的轨迹会导致明显的性能下降，且需要许多后续 RL 步才能恢复。仅从最终 RL 检查点生成轨迹，在恢复 RL 后性能较弱。一个可能的解释是：在一系列强检查点上收集的轨迹比从单一最终策略生成的轨迹具有更高的多样性，从而在 RL 中带来更好的探索。复用 RL 期间已生成的轨迹还能避免重采样所需的额外计算和时间。

• 在固定 token 预算下，提高提示词多样性的价值高于增加每个提示词的轨迹数量。我们还发现简单随机采样优于若干有偏选择策略，包括最短轨迹采样以及与 Yang et al.~(2025) 探索的类似启发式过滤方法。

• 在自蒸馏期间，模型可能遗忘中训练学到的长上下文行为，尤其是当推理轨迹来自以较短最大长度训练的早期 RL 运行时。为缓解这一点，我们在长度扩展之前，将中训练数据与推理轨迹混合进行自蒸馏。

\subsection{3.1.5 超参数}\label{hyperparameters-1}

下面我们报告三个领域特定模型的主要 RL 攀登以及自蒸馏运行中通常使用的超参数。

RL 攀登。我们使用 AdamW 训练，其中 \(\beta _ { 1 } = \beta _ { 2 } = 0 . 9 5 , \epsilon = 1 0 ^ { - 1 5 }\)，且无权重衰减。我们使用恒定学习率 \(\eta = 1 0 ^ { - 6 }\)，无 warmup 或衰减。打包后的全局批大小为 7040，未打包序列总数上限为 12000。由于我们的 RL 技术栈完全异步，随着生成长度增加，rollouts 的延迟会上升，而训练步耗时不会成比例增加（因为存在大量短生成）。这意味着较长的 rollouts——通常是难题——的 off-policy 程度会上升。为应对这一点，我们选择在较高长度下将学习率降至 $9 \times 1 0 ^ { - 7 }$，以降低 off-policy 程度并提高稳定性。最大生成长度为 128k，但我们先以 8k、16k、32k 和 64k 训练。

对于带自适应熵控制的 GRPO，我们设置 \(\epsilon = 0 . 6 , k _ { \mathrm { m a x } } = 2 . 5\)、步长 \(\delta \ : = \ : 0 . 2 5\) 以及目标熵 \(H ^ { \star } = 0 . 3\)。对于外层裁剪，我们设置 \(r _ { \operatorname* { m a x } } = 5 0\)，并让 \(r _ { \operatorname* { m i n } } = 0\) 不受约束。

对于语言一致性奖励，我们设置 \(w _ { \mathrm { l a n g } } = 0 . 5\)，逐词惩罚 \(\alpha = 0 . 0 0 5\)。对于长度惩罚，我们在 64k 长度扩展阶段之前固定 \(w _ { \mathrm { l e n } } = 0 . 2 5\)。由于该惩罚由 \(\ell _ { \mathrm { m a x } }\) 归一化，有效惩罚会随最大长度增加而自然减弱，使模型能够在更大的 token 预算下探索更长的推理轨迹。我们在 128k 扩展阶段移除长度惩罚（\(w _ { \mathrm { l e n } } = 0\)）。

对于问题采样，我们设置共 \(G = 1 2 8\) 次 rollouts，其中 \(G _ { \mathrm { e a r l y } } = 1 6\)，\([ \rho _ { \mathrm { m i n } } ^ { \mathrm { e a r l y } } , \rho _ { \mathrm { m a x } } ^ { \mathrm { e a r l y } } ] = [ 0 . 0 5 , 0 . 8 ]\)。通过率过滤使用 \([ \rho _ { \operatorname* { m i n } } , \rho _ { \operatorname* { m a x } } ] = [ 0 . 1 , 0 . 8 ]\)。对于 top-p 采样，我们使用 $p = 0 . 9 7$。更大的核取值能改善探索，但会增加传输采样掩码的基础设施开销；我们发现 $p = 0 . 9 7$ 在探索与训练效率之间取得了良好平衡。

我们在两次推理模型更新之间执行 5 个梯度步，并丢弃任何其生成策略超过 8 次推理更新（即落后 40 个梯度步）的 rollout。这在略微放宽 on-policy 新鲜度的同时显著提高了吞吐。我们使用 dropless MoE 训练，并将全局 MoE 负载均衡系数设为 \(1 \cdot 1 0 ^ { - 5 }\)。

自蒸馏。对于自蒸馏 SFT，我们使用打包序列，全局批大小为 2,048，序列长度为 128k token。优化使用 AdamW，权重衰减 0.001，学习率采用余弦调度。最大学习率设为 \(1 . 7 \cdot 1 0 ^ { - 5 }\)，最小学习率设为 \(5 . 2 \cdot 1 0 ^ { - 6 }\)，学习率 warmup 比例为 2\%。



\pandocbounded{\includegraphics[keepaspectratio,alt={image}]{images/131b7346ca39b1019e838c1fccccbffb286f0c2f2acd0e78bb3d2cd9f8b0d84f.jpg}}

图 16. 我们的 STEM Mix 数据集按原始题目格式（左）和学科分类体系（右）的分布，其中 Other 汇总了各自占比低于数据集 0.5\% 的主题。多项选择题和证明题在摄入期间被转换为开放式形式（转换不可行的予以剔除）；我们保留一小部分多项选择题，使模型对该格式保持熟悉。

\pandocbounded{\includegraphics[keepaspectratio,alt={image}]{images/5e47e21c5d20ed6ffed156f4c69d57f835277136ef8f7cd75fe0b69a3f9e73be.jpg}}

在自蒸馏期间，我们仔细调优两个超参数：dropout（Srivastava et al., 2014）和 MoE 负载均衡损失系数。我们使用相对较高的 0.15 的 dropout 率，这会提高熵并有助于防止模型坍缩，从而提升后续 RL 阶段的性能。由于领域特定 RL 中的数据分布远比预训练和中训练使用的混合数据窄，RL 期间的专家选择可能变得高度不均衡。然而，在 RL 中使用较大的负载均衡系数也会损害稳定的性能提升。为解决这一问题，我们在自蒸馏期间使用相对较大的负载均衡系数 \(1 \cdot 1 0 ^ { - 2 }\)，而在 RL 期间使用小得多的系数 \(1 \cdot 1 0 ^ { - 5 }\)。由于自蒸馏的上下文来自 RL 运行本身，因此分布相似，自蒸馏期间诱导的专家均衡效应在 RL 期间得以保留。
